\section{Shifted boundary resultants and the complete unmatched depth}
\label{srds-section}

Let \(n>324\) be prime, \(B=n^4\), and
\(\mathcal O=\mathbb Z[\zeta]\), where
\(\zeta^2-\zeta+1=0\) and \(\operatorname{Im}\zeta>0\).
For \(z=A+C\zeta\), put \(P(z)=AC(A+C)\), and define
\[
 P_n(a)=P((a+\zeta)^n),\qquad
 T_n(X)=P_n(3X),\qquad F_n(X)=T_n(X)/3,\qquad d=3n-1.
\]
Section~\ref{cg19-section} proves that \(F_n\) is a primitive
integer polynomial of degree \(d\), with
\[
 \|F_n\|_2\le\|F_n\|_1\le64^n.
\]
For every integer \(B\le k<2B\), write \(w_k=3k+\zeta\) and
\(T_k=T_n(k)\). Its three arms are positive and pairwise coprime;
every boundary prime is a unit for the actual input norm
\(N(w_k)=9k^2+3k+1\).

Fix a nonzero integer \(h\) with \(|h|<B\), and set
\[
 I_h=\{k\in\mathbb Z:B\le k<2B,\ B\le k+h<2B\}.
\]
This is a nonempty set of \(B-|h|\) actual pairs. The result below
controls their common prime-power depths simultaneously even when
different primes use different witness pairs. It retains all
unmatched depth in the signed accounting.

\begin{proposition}[Every nonzero integer translate is coprime]
\label{srds-coprime}
For every nonzero integer \(h\), without the additional restriction
\(|h|<B\), one has
\[
 \gcd_{\mathbb Q[X]}(F_n(X),F_n(X+h))=1.
\]
\end{proposition}
\begin{proof}
Work first in the \(a\)-coordinate. Put
\(\xi=\zeta^2\), \(K=\mathbb Q(\zeta)\), and
\(L=\mathbb Q(\zeta,\eta)\), where \(\eta\) is a primitive
\(n\)-th root of unity. The group
\(\operatorname{Gal}(L/K)\) has order \(n-1\) and sends \(\eta\)
to all \(\eta^j\), \(1\le j<n\). Indeed the cyclotomic field of
conductor \(3n\) has degree \(2(n-1)\), and CRT permits every
nonzero residue modulo \(n\) while fixing the residue modulo three.

Write
\((a+\zeta)^n=A_n(a)+C_n(a)\zeta\).
A root of \(P_n\) is neither \(-\zeta\) nor \(-\bar\zeta\):
at either exceptional value one power embedding is zero and the
other nonzero, whereas any one of the three arm equations
\(A_n=0\), \(C_n=0\), \(A_n+C_n=0\) would force both embeddings
to vanish. The ratio
\[
 \alpha=\frac{a+\zeta}{a+\bar\zeta}
\]
is therefore defined and nonzero, and the exact arm conditions are
\begin{equation}
 \alpha^n=
 \begin{cases}
  1,&C_n=0,\\
  \xi,&A_n=0,\\
  \xi^2,&A_n+C_n=0.
 \end{cases}
 \label{srds-arm-phases}
\end{equation}
Conversely each equation implies its indicated arm equation by
subtracting the two power embeddings. The value \(\alpha=1\)
is impossible, and the inverse transformation is
\begin{equation}
 a=\frac{\zeta-\alpha\bar\zeta}{\alpha-1}.
 \label{srds-ratio-inverse}
\end{equation}
Thus the roots are exactly the \(3n-1\) distinct images of
\(\mu_{3n}\setminus\{1\}\). Their distinctness and the known degree
also account for all multiplicities.

The two elements in \(\mu_3\setminus\{1\}\) yield \(a=0,-1\).
All other elements form the three full \(K\)-Galois orbits
\[
 \{\xi^j\eta^t:1\le t<n\},\qquad j=0,1,2,
\]
each of cardinality \(n-1\). The inverse
\eqref{srds-ratio-inverse} preserves these orbit sizes.
Since \(n\) is prime to three, one orbit lies in each arm after
its possible rational root has been removed. This gives all
\(3(n-1)+2\) roots.

For every root \(a\), including a rational root, define
\(\mu(a)=\operatorname{Tr}_{L/K}(a)/(n-1)\).
The leading and next coefficients of the three arm polynomials
give their total root sums:
\[
 \begin{array}{c|c|c}
 \text{arm}&\text{leading terms}&\text{total root sum}\\ \hline
 A_n&a^n+0a^{n-1}&0\\
 C_n&na^{n-1}+\frac{n(n-1)}2a^{n-2}&-(n-1)/2\\
 A_n+C_n&a^n+na^{n-1}&-n
 \end{array}
\]
When \(n\equiv1\pmod6\), \(A_n\) has rational root zero and
\(A_n+C_n\) has rational root \(-1\). When \(n\equiv5\pmod6\)
these rational roots are interchanged. Consequently the three
nonrational orbit means are
\begin{equation}
 \begin{array}{c|ccc}
 n\bmod6&\multicolumn{3}{c}{\text{nonrational orbit means}}\\ \hline
 1&0&-1/2&-1\\
 5&1/(n-1)&-1/2&-n/(n-1).
 \end{array}
 \label{srds-trace-means}
\end{equation}
The rational roots have means zero and \(-1\). The diameter of
the complete set of possible means is at most
\((n+1)/(n-1)<3\).

If \(F_n(X)\) and \(F_n(X+h)\) had a common root \(x\), the
two roots \(a=3x\) and \(b=3x+3h\) of \(P_n\) would both lie in
the same field \(L\). Linearity of its full normalized trace would
give \(\mu(b)-\mu(a)=3h\), whose absolute value is at least three.
This contradicts \eqref{srds-trace-means}. No assumption that the
two roots have the same minimal polynomial was used.
\end{proof}

\begin{proposition}[The actual common-depth lcm]
\label{srds-lcm}
For \(0<|h|<B\), define the positive integers
\[
 C_F(h)=\mathop{\operatorname{lcm}}_{k\in I_h}
       \gcd(F_n(k),F_n(k+h)),\qquad
 C_T(h)=\mathop{\operatorname{lcm}}_{k\in I_h}
       \gcd(T_k,T_{k+h}).
\]
Then
\begin{equation}
 \begin{split}
 C_F(h)&\mid\operatorname{Res}(F_n(X),F_n(X+h)),\\
 C_T(h)&=3C_F(h).
 \end{split}
 \label{srds-lcm-div}
\end{equation}
In particular, with
\begin{equation}
 H_n(h)=2dn\log64+d^2\log(1+|h|),
 \label{srds-height-cost}
\end{equation}
one has
\begin{equation}
 \log C_T(h)\le H_n(h)+\log3.
 \label{srds-lcm-height}
\end{equation}
For each fixed nonzero integer \(h\), this bound is \(O_h(n^2)\).
\end{proposition}
\begin{proof}
For each prime \(p\), let
\[
 b_p=\max_{k\in I_h}
       \min\{v_p(F_n(k)),v_p(F_n(k+h))\}.
\]
For every positive \(b_p\), choose an actual pair attaining it.
Different primes may require different pairs. Both polynomials
vanish at the chosen owner modulo \(p^{b_p}\), and
Proposition~\ref{srds-coprime} proves their coprimality.

For clarity, the integer resultant argument applies to full prime
powers without monicity assumptions. The Sylvester matrix is the
coefficient matrix of
\((U,V)\mapsto UF+VG\) with the usual complementary degree
bounds. Its adjugate applied to the constant coefficient vector
gives integer polynomials with
\[
 UF+VG=\operatorname{Res}(F,G),
\]
up to a sign convention. Evaluation at the chosen owner proves
\(p^{b_p}\) divides the nonzero resultant. Multiplication over
the distinct primes proves the first assertion in
\eqref{srds-lcm-div}.

For every pair, \(\gcd(3F_n(k),3F_n(k+h))\) is three times the
corresponding \(F_n\)-gcd. The lcm of a nonempty finite list
commutes with this common scaling, proving the second assertion.
Thus at three the maximum common \(T\)-depth is exactly one more
than the maximum common \(F\)-depth; it is not assumed to be one.

Translation satisfies
\[
 \|F_n(X+h)\|_2\le\|F_n(X+h)\|_1
       \le\|F_n\|_1(1+|h|)^d,
\]
because the coefficient norm of \((X+h)^j\) is
\((1+|h|)^j\). There are \(d\) shifted coefficient rows of each
polynomial in the Sylvester matrix. Hadamard's inequality therefore
gives
\[
 \begin{split}
 \log|\operatorname{Res}(F_n,F_n(\,\cdot+h))|
 &\le d\log\|F_n\|_2+d\log\|F_n(\,\cdot+h)\|_2\\
 &\le2dn\log64+d^2\log(1+|h|).
 \end{split}
\]
Combining this with \eqref{srds-lcm-div} proves the claim.
\end{proof}

The second polynomial has no packet modulus hidden in its degree
or coefficient bound. Its degree can be much smaller than the
number of distinct witness owners. This gives one simultaneous
integer bound for the actual common-depth part, without any
multiplicative-independence condition. It does not control excess
depth at one owner beyond the depth at its partner.

\begin{proposition}[The full signed residual]
\label{srds-signed}
Let \(\mathcal P\) be a finite set of distinct primes, each with
an assigned owner \(k_p\in I_h\) and full depth
\(e_p=v_p(T_{k_p})\ge1\). Set
\[
 b_p=\min\{e_p,v_p(T_{k_p+h})\},\qquad u_p=e_p-b_p,\qquad
 \mathcal P_c=\{p\in\mathcal P:b_p>0\},\qquad
 R_c=\prod_{p\in\mathcal P_c}p.
\]
The original signed packet has the exact identity
\begin{equation}
 \begin{split}
 J_{\mathcal P}:=\sum_{p\in\mathcal P}(e_p-3)\log p
   &=\sum_{p\in\mathcal P_c}(b_p-3)\log p
       +\sum_{p\in\mathcal P_c}u_p\log p\\
   &\quad+\sum_{p\in\mathcal P\setminus\mathcal P_c}
                    (e_p-3)\log p.
 \end{split}
 \label{srds-exact-signed}
\end{equation}
If \(\delta_3\) indicates \(3\in\mathcal P_c\), its proved bound is
\begin{equation}
 \begin{split}
 J_{\mathcal P}
 &\le H_n(h)+\delta_3\log3-3\log R_c\\
 &\quad+\sum_{p\in\mathcal P_c}u_p\log p
       +\sum_{p\in\mathcal P\setminus\mathcal P_c}
                    (e_p-3)\log p.
 \end{split}
 \label{srds-unmatched-bound}
\end{equation}
\end{proposition}
\begin{proof}
Identity \eqref{srds-exact-signed} is termwise addition. Apply the
resultant certificate to precisely the selected common packet,
reducing its depth at three by one and omitting a zero depth.
The \(T\)-modulus is three times the adjusted \(F\)-modulus when
three is present. Restoring the original radical, exactly as in
\eqref{cg19-full-signed}, gives
\[
 \sum_{p\in\mathcal P_c}(b_p-3)\log p
       \le H_n(h)+\delta_3\log3-3\log R_c.
\]
Substitution proves \eqref{srds-unmatched-bound}. If the common
packet is empty, its cost and radical logarithm are zero and the
same inequality remains valid. All depth-one and depth-two terms
remain on their actual supports with their negative signs.
\end{proof}

In particular the complete lcm has
\[
 J(C_T(h))\le H_n(h)+\log3-3\log\operatorname{rad}(C_T(h)),
 \qquad J(M)=\log M-3\log\operatorname{rad}(M).
\]
The larger radical of this complete lcm cannot replace the
radical of a selected subpacket. Signed sums are not monotone
under the addition of depth-one or depth-two labels.

\paragraph{Several allowed shifts and their cost.}
Let \(\mathcal H\) be a finite set of nonzero integer shifts
with \(|h|<B\). Allow each distinct prime its own selected
shift \(h_p\in\mathcal H\) and actual pair in \(I_{h_p}\).
Partitioning the primes by that shift gives
\begin{equation}
 \begin{split}
 \log M_{\rm common}
     &\le\sum_{h\in\mathcal H}H_n(h)+\delta_3\log3,\\
 J_{\rm common}
     &\le\sum_{h\in\mathcal H}H_n(h)+\delta_3\log3
                 -3\log\operatorname{rad}(M_{\rm common}).
 \end{split}
 \label{srds-multishift}
\end{equation}
Only one \(\log3\) is restored, because three is a single assigned
prime label. Every unmatched term in \eqref{srds-exact-signed}
must still be added.

For the actual full block scale
\[
 W_n=\sum_{B\le k<2B}n\log|3k+\zeta|
       \ge Bn\log(3B),
\]
a fixed finite set \(\mathcal H\) has
\(\sum_hH_n(h)/W_n\longrightarrow0\).
More generally the same conclusion follows, for nonempty
\(\mathcal H\), if
\[
 \frac{|\mathcal H|\,n\log(1+\max_{h\in\mathcal H}|h|)}
      {B\log B}\longrightarrow0.
\]
Summing the proved bounds over all \(0<|h|<B\) instead costs
\(O(Bn^2\log B)\), not \(o(Bn\log B)\).
Neither a covering by few shifts nor control of all unmatched
depth follows from this calculation.

\begin{proposition}[Different positive labels in the original block]
\label{srds-actual-example}
For every prime \(n>324\), the original block contains two
distinct pairs with common exact depths four at seven and five
at thirteen, respectively; neither prime hits the other pair.
\end{proposition}
\begin{proof}
Use the fixed integers
\[
 L_0=7^4\,13^5=891474493,\qquad
 h=(L_0-1)/3=297158164.
\]
Let \(k_7\) be the first integer at least \(B\) satisfying
\[
 k_7\equiv7^4\pmod{7^5},\qquad k_7\equiv1\pmod{13},
\]
and define \(k_{13}\) by
\[
 k_{13}\equiv13^5\pmod{13^6},\qquad k_{13}\equiv1\pmod7.
\]
The respective CRT steps are \(218491\) and \(33787663\).
Since
\[
 h+33787663<331000000<324^4<B,
\]
all four indices \(k_7,k_7+h,k_{13},k_{13}+h\) lie in the
original block.

The integral expansion \(P_n(a)=na+O(a^2)\) was proved in
Section~\ref{cg19-section}. Also
\[
 P_n(-1-a)=P_n(a),\qquad P_n'(-1)=-n.
\]
Indeed \(-1-a+\zeta=-\overline{a+\zeta}\), \(n\) is odd,
and \(P(-\bar z)=P(z)\). At \(k_q\) the expansion at zero
gives exactly the depth \(e\) for
\((q,e)=(7,4),(13,5)\), since \(n\) is a \(q\)-unit.
At the shifted index,
\[
 3(k_q+h)+1=3k_q+L_0.
\]
After division by \(q^e\), its residue is respectively
\[
 3+13^5\equiv2\pmod7,\qquad
 3+7^4\equiv12\pmod{13}.
\]
Both are units, so the expansion at \(-1\) gives the same exact
depth at the second owner. Primitivity and the norm-unit
conditions follow from the actual original-block setup.

For \(p=7,13\), \(\gcd(n,p-1)=1\). At an input of nonzero
norm, the split ratio group has order \(p-1\).
Equation \eqref{srds-arm-phases} then forces a boundary hit
to have \(\alpha\in\mu_3\), hence \(a=0\) or \(-1\)
modulo \(p\). At norm zero, primitivity and the norm identity
exclude a boundary hit instead. At the other pair the two
input values \(a=3k\) are \(3\) and \(2\) modulo \(p\), since
\(k\equiv1\) and \(h\equiv-1/3\pmod p\). Neither is zero or
\(-1\), proving absence of that prime throughout the other pair.
The four owners are distinct: modulo seven the first pair has
indices \(0,2\), and the second \(1,3\).
\end{proof}

The selected packet consequently has
\[
 M_{\rm selected}=7^4\,13^5,\qquad
 J_{\rm selected}=\log7+2\log13>0.
\]
No factorization of the remaining boundary is assumed. This is
an actual original-block nonemptiness witness, with different
prime owners. Its marked primes are fixed and small; it is not
a far-tail example or a claim of membership in the independent
domain.

\begin{proposition}[Assigned matched costs and their exact quantifier]
\label{srds-assigned-markov}
Fix \(n,B,h\) as above and put \(t_k=n\log|3k+\zeta|\),
\(C_n(h)=H_n(h)+\log3\). For any finite set \(\mathcal P\)
of distinct primes, choose one actual pair per prime, with
\[
 b_p=\min\{v_p(T_{k_p}),v_p(T_{k_p+h})\}\ge1.
\]
Define its assigned positive matched cost
\[
 L_h(k)=\sum_{\substack{p\in\mathcal P\\k_p=k}}
                          (b_p-3)_+\log p,
 \qquad x_+=\max(x,0).
\]
For every real \(\epsilon>0\),
\begin{equation}
 \begin{split}
 \sum_{k\in I_h}L_h(k)&\le C_n(h),\\
 \#\{k\in I_h:L_h(k)>\epsilon t_k\}
       &<\frac{C_n(h)}{\epsilon n\log(3B)}.
 \end{split}
 \label{srds-assigned-bound}
\end{equation}
\end{proposition}
\begin{proof}
Each selected common depth is at most its corresponding exponent
in \(C_T(h)\). Distinctness of the assigned prime labels gives
\[
 \sum_k L_h(k)=\sum_{p\in\mathcal P}(b_p-3)_+\log p
    \le\sum_{p\in\mathcal P}b_p\log p
    \le\log C_T(h)\le C_n(h).
\]
Every exceptional owner's cost exceeds
\(\epsilon t_k\ge\epsilon n\log(3B)\). Summing over the
exceptional set proves the strict cardinality bound when it is
nonempty; when it is empty, strictness holds because the displayed
right side is positive.
\end{proof}

For fixed \(h\), the exception bound is
\(O_h(n/(\epsilon\log n))\). Taking \(\epsilon_n=1/n\)
gives \(O_h(n^2/\log n)=o(B)\) exceptions and cost at most
\(t_k/n\) on every other assigned owner. The same finite
argument for multiple shifts uses their total cost in
\eqref{srds-multishift}; an empty shift set has zero cost
and zero exceptions, with a non-strict zero bound.

The quantifier is: for every specified prime/pair assignment
there exists its corresponding exceptional set. It is not
one common good set for all assignments. In this proposition
a prime is assigned once; its common depth is not summed over
every pair where it occurs. The positive-cost estimate also
does not replace the complete signed identity
\eqref{srds-exact-signed}.

\paragraph{The precise remaining second-polynomial requirement.}
For any assigned full \(T\)-packet, let \(R_T\) be its original
radical and let \(\delta_3\) indicate whether three is present.
A nonzero \(G\in\mathbb Z[X]\) coprime to \(F_n\) must satisfy
the adjusted \(F_n\)-packet: the depth at three is \(e_3-1\)
when positive and is omitted when zero; all other depths are
unchanged. Define
\[
 D_n(G)=\deg(G)\log\|F_n\|_2+d\log\|G\|_2.
\]
The corrected full signed certificate is
\begin{equation}
 J_{\rm packet}\le D_n(G)+\delta_3\log3-3\log R_T.
 \label{srds-second-certificate}
\end{equation}
It follows from the integer Sylvester bound and, for a nonzero
constant \(G=c\), the direct case
\(\operatorname{Res}(F_n,c)=c^d\).
If a larger actual signed sum has the exact decomposition
\(J_{\rm total}=J_{\rm packet}+J_{\rm remainder}\), a sufficient
certificate at a positive scale \(W\) and tolerance \(\epsilon>0\)
is
\begin{equation}
 D_n(G)\le3\log R_T-\delta_3\log3
                   -J_{\rm remainder}+\epsilon W.
 \label{srds-required-cost}
\end{equation}
For specified degree \(a\), this is equivalently
\[
 \log\|G\|_2\le
 \frac{3\log R_T-\delta_3\log3-J_{\rm remainder}
             +\epsilon W-a\log\|F_n\|_2}{d}.
\]
Replacing \(\log\|F_n\|_2\) by \(n\log64\) gives a stronger
sufficient condition. A negative right side cannot be attained
by a nonzero integer coefficient vector.

The radical and remainder keep all original depth-one and
depth-two credits. Criterion \eqref{srds-required-cost} does
not construct \(G\). The automatic product of owner factors has
resultant equal to the product of owner values and recovers their
old height bound. Even degree reduction alone does not suffice:
the actual CG packet also permits \(G=625(X-B)\), whose exact
resultant \(625^dF_n(B)\) still pays the large original owner value.

\paragraph{Coverage and the independent-domain singleton range.}
Write \(L=\log(8B)\) and \(r=\lceil\sqrt{n/2}\rceil\).
The previously proved unit-signed theorem says that in its
multiplicatively independent domain \(\mathcal I\), a prime
\(q>8B\) with \(\log q\ge rL/2\) has at most one root of
depth at least four. Thus, when both endpoints of a pair lie
in \(\mathcal I\), its common positive depth at such a prime
is already zero. Under the stronger inequality
\(\log q\ge2rL\), that domain has at most one positive-depth
root, so it contains no distinct common-support pair at all.

The shifted resultant permits a partner in the complete original
block, including outside \(\mathcal I\). It proves neither that
such a partner exists nor that it has comparable depth or uses
one of a small set of shifts. In
\eqref{srds-exact-signed}, all differences \(e_p-b_p\) and
all unpaired labels are still present. There is no lower bound
on the proportion of original positive depth contained in the
matched portion.

The \(O_h(n^2)\) fixed-shift cost is small on the full block
scale \(Bn\log B\) but exceeds the individual scale of order
\(n\log B\). Proposition~\ref{srds-assigned-markov} is a finite
exceptional-set consequence for a specified assignment, not
a bound for every root. All-shift summation removes its saving.
Exceptional roots, the dependent-root complement, and the
reduction covering arbitrary primitive ABC triples remain
separate obligations.

\paragraph{Evidence and scope.}
The ordinary root-orbit, trace, integer resultant and signed
proofs above were independently reviewed in full. The exact
script \texttt{next\_replay\_shifted\_\allowbreak pairs.py}
checks twelve literal binary Eisenstein powers for
\(n=331,337,347\), the block conditions, primitivity, norm
identities, exact marked depths and cross-prime absence.
Its canonical output SHA256, concatenating the two lines, is
\begin{center}
\small\ttfamily
980fc36976bda4a192a6caa0adae886547\\
613aade9b97aabcec310c14e12265b
\end{center}
The author and root independently ran its check successfully.
That finite replay does not prove the infinite root-orbit or
resultant assertions, which have their ordinary proofs above.
This section claims no new Lean theorem and no full ABC proof.
Later simple-root or differential-certificate candidates are
not included in these statements.
